\scanpage{102}\mspage{98}
\[
 f\leq g\quad\overset{\text{\unsure{dfn|defn|def}}}{\Longleftrightarrow}\quad
 \forall x\in\Sigma,\qquad f(x)\leq g(x)\;.
\]

\unsure{Je|J} dis que pour cette relation d'ordre dans $\operatorname{Corr}(\Sigma,\Sigma')$, les $\mathrm{Sup}$ quelconques (donc aussi les $\mathrm{Inf}$ quelconques) existent, et sont \unsure{données|donnés} par
\[
 \bigl(\mathop{\mathrm{Sup}}\limits_i \unsure{f^{\cdot}|f_i|f'}\bigr)(x)
 =\mathop{\mathrm{Sup}}\limits_i \unsure{f^{\prime}|f^{\cdot}|f_i}(x)\qquad\forall \unsure{x|\cdot}\in\Sigma.
\]
En effet, l'associativité des $\mathrm{Sup}$ montre que l'application du second \unsure{mb|=b} commute bien aux $\mathrm{Sup}$ \unsure{quelc.|quelc}, donc est une corr. de spatios.

Mais la structure $C$, ou la structure de \unsure{disjonction|disponction}, n'est pas claire.

Dans le cas où $E$ est un \unsure{disjoncteur|disponcteur} trivial, \deleted{on a \unsure{une}} et où les corr. $E\unsure{\mbox{---}|\longrightarrow}\Sigma'$ sont \inserted{\unsure{données par}} les applications quelconques $\deleted{\Sigma}E\unsure{\mbox{---}|\longrightarrow}\Sigma'$, donc
\[
 \operatorname{Corr}(\deleted{E}\Sigma_E,\Sigma')\simeq{\Sigma'}^{E},
\]
les structures de disjonction et $C$ peuvent \unsure{de|se} décrire ``terme par \unsure{terme|termes}''
\[
 (\unsure{\complement|C} f)(x)=\complement f(x)\qquad\forall x\in E
\]
% archive sheet number: 101
